\documentclass[12pt,a4paper]{article}
\usepackage[utf8]{inputenc}
\usepackage[spanish]{babel}
\usepackage{geometry}
\usepackage{fancyhdr}
\usepackage{enumitem}
\usepackage{xcolor}
\usepackage{listings}
\usepackage{graphicx}
\usepackage{hyperref}
\usepackage{tabularx}
\usepackage{array}

\geometry{margin=2cm}
\pagestyle{fancy}
\fancyhf{}
\lhead{Curso de Programación en Python}
\rhead{Proyecto Final}
\cfoot{\thepage}

% Configuración para código Python
\lstset{
    language=Python,
    basicstyle=\ttfamily\small,
    keywordstyle=\color{blue},
    stringstyle=\color{red},
    commentstyle=\color{green!60!black},
    showstringspaces=false,
    breaklines=true,
    frame=single,
    numbers=left,
    numberstyle=\tiny\color{gray}
}

\title{\textbf{Proyecto Final del Curso}\\
\large Sistema de Gestión de Biblioteca Personal con Interfaz Gráfica}
\author{Curso de Programación en Python}
\date{Tiempo de entrega: 2 semanas}

\begin{document}

\maketitle

\section{Introducción}

¡Felicidades por llegar al proyecto final! Este proyecto integra \textbf{todos los conceptos} que has aprendido durante el curso. Vas a crear una aplicación completa y profesional: un \textbf{Sistema de Gestión de Biblioteca Personal} con interfaz gráfica usando \texttt{tkinter}.

Este proyecto no solo demuestra tus habilidades técnicas, sino también tu capacidad para integrar múltiples conceptos, diseñar software robusto y crear aplicaciones con valor práctico real.

\section{Descripción General del Proyecto}

Crearás un sistema completo para gestionar una biblioteca personal de libros. La aplicación permitirá:

\begin{itemize}[leftmargin=*]
    \item Agregar, buscar, editar y eliminar libros
    \item Categorizar libros por género, autor, año, etc.
    \item Registrar préstamos a amigos/familiares
    \item Generar estadísticas sobre la colección
    \item Importar/exportar datos desde/hacia archivos CSV
    \item Validar datos de entrada
    \item Buscar libros usando expresiones regulares
    \item Guardar y recuperar datos automáticamente
    \item Interfaz gráfica intuitiva y profesional
\end{itemize}

\section{Objetivos de Aprendizaje}

Al completar este proyecto, demostrarás dominio de:

\subsection{Conceptos Fundamentales (Semanas 0-2)}
\begin{itemize}[leftmargin=*]
    \item Variables y tipos de datos
    \item Funciones y parámetros
    \item Condicionales complejos
    \item Bucles y validación iterativa
    \item Formateo de strings con f-strings
\end{itemize}

\subsection{Conceptos Intermedios (Semanas 3-5)}
\begin{itemize}[leftmargin=*]
    \item Manejo robusto de excepciones
    \item Uso de librerías estándar y externas (\texttt{datetime}, \texttt{csv}, \texttt{re}, \texttt{statistics}, etc.)
    \item Pruebas unitarias con \texttt{pytest}
    \item Validación y sanitización de datos
\end{itemize}

\subsection{Conceptos Avanzados (Semanas 6-9)}
\begin{itemize}[leftmargin=*]
    \item Lectura y escritura de archivos
    \item Trabajo con CSV
    \item Expresiones regulares para búsqueda y validación
    \item Programación orientada a objetos (clases, herencia, encapsulación)
    \item Comprehensions (list, dict)
    \item Funciones lambda
    \item Decoradores
    \item Interfaz gráfica con \texttt{tkinter}
\end{itemize}

\section{Requisitos Funcionales}

\subsection{Módulo 1: Clase Libro (POO)}

Crea una clase \texttt{Libro} con los siguientes atributos y métodos:

\begin{lstlisting}[caption=Estructura básica de la clase Libro]
class Libro:
    def __init__(self, isbn, titulo, autor, genero, anio_publicacion):
        # Atributos privados con validacion
        self._isbn = isbn
        self._titulo = titulo
        self._autor = autor
        self._genero = genero
        self._anio_publicacion = anio_publicacion
        self._prestado = False
        self._prestado_a = None
        self._fecha_prestamo = None
    
    # Properties con @property
    @property
    def isbn(self):
        return self._isbn
    
    # Metodos de validacion
    @staticmethod
    def validar_isbn(isbn):
        # Usar regex para validar formato ISBN
        pass
    
    # Metodos especiales
    def __str__(self):
        pass
    
    def __repr__(self):
        pass
    
    # Metodos de negocio
    def prestar(self, nombre_persona):
        pass
    
    def devolver(self):
        pass
\end{lstlisting}

\textbf{Requisitos de la clase:}
\begin{itemize}[leftmargin=*]
    \item Validar ISBN usando expresiones regulares (formato ISBN-10 o ISBN-13)
    \item Validar que el año de publicación sea razonable (entre 1000 y año actual)
    \item Usar \texttt{@property} para todos los atributos principales
    \item Método \texttt{\_\_str\_\_()} que retorne representación legible
    \item Método \texttt{\_\_repr\_\_()} para debugging
    \item Método para prestar libro (registra fecha usando \texttt{datetime})
    \item Método para devolver libro
    \item Al menos un método de clase (\texttt{@classmethod}) y uno estático (\texttt{@staticmethod})
\end{itemize}

\subsection{Módulo 2: Clase Biblioteca (POO y Gestión)}

Crea una clase \texttt{Biblioteca} que gestione la colección:

\begin{lstlisting}[caption=Estructura de la clase Biblioteca]
class Biblioteca:
    def __init__(self, nombre="Mi Biblioteca"):
        self._nombre = nombre
        self._libros = {}  # Diccionario: ISBN -> Libro
        self._archivo_datos = "biblioteca.csv"
        self.cargar_desde_archivo()
    
    def agregar_libro(self, libro):
        # Agregar con validacion de duplicados
        pass
    
    def buscar_por_titulo(self, patron):
        # Busqueda con regex
        pass
    
    def buscar_por_autor(self, patron):
        # Busqueda con regex
        pass
    
    def eliminar_libro(self, isbn):
        # Eliminar con confirmacion
        pass
    
    def listar_libros(self, filtro=None):
        # Retorna lista filtrada
        pass
    
    def obtener_estadisticas(self):
        # Usar statistics para calcular promedios, etc.
        pass
    
    def guardar_a_archivo(self):
        # Guardar en CSV
        pass
    
    def cargar_desde_archivo(self):
        # Cargar desde CSV con manejo de excepciones
        pass
    
    def exportar_csv(self, archivo):
        # Exportar a CSV personalizado
        pass
    
    def importar_csv(self, archivo):
        # Importar desde CSV con validacion
        pass
\end{lstlisting}

\textbf{Requisitos de la clase:}
\begin{itemize}[leftmargin=*]
    \item Usar diccionario para almacenar libros (clave: ISBN)
    \item Búsqueda con expresiones regulares (insensible a mayúsculas)
    \item Guardar automáticamente después de cada cambio
    \item Manejo robusto de excepciones en I/O de archivos
    \item Método que use \texttt{statistics} para calcular estadísticas (año promedio, etc.)
    \item Usar comprehensions para filtrar y transformar datos
    \item Funciones lambda en operaciones de ordenamiento y filtrado
\end{itemize}

\subsection{Módulo 3: Interfaz Gráfica (Tkinter)}

Crea una interfaz gráfica completa con las siguientes ventanas/secciones:

\subsubsection{Ventana Principal}
\begin{itemize}[leftmargin=*]
    \item Menú superior con opciones: Archivo, Gestión, Estadísticas, Ayuda
    \item Panel lateral con botones para funciones principales
    \item Área central con tabla/lista de libros
    \item Barra de estado inferior
\end{itemize}

\subsubsection{Ventana de Agregar/Editar Libro}
\begin{itemize}[leftmargin=*]
    \item Campos de entrada para todos los atributos del libro
    \item Validación en tiempo real
    \item Mensajes de error claros
    \item Botones Guardar y Cancelar
\end{itemize}

\subsubsection{Ventana de Búsqueda}
\begin{itemize}[leftmargin=*]
    \item Campo de búsqueda con opciones (por título, autor, género)
    \item Checkbox para usar expresiones regulares
    \item Lista de resultados
    \item Botones para acciones sobre resultados
\end{itemize}

\subsubsection{Ventana de Estadísticas}
\begin{itemize}[leftmargin=*]
    \item Total de libros en la colección
    \item Libros por género (gráfico o tabla)
    \item Año promedio de publicación
    \item Libros actualmente prestados
    \item Top 5 autores con más libros
\end{itemize}

\subsubsection{Ventana de Préstamos}
\begin{itemize}[leftmargin=*]
    \item Registrar nuevo préstamo
    \item Ver préstamos activos
    \item Devolver libro
    \item Historial de préstamos
\end{itemize}

\subsection{Módulo 4: Validaciones y Excepciones}

Implementa manejo robusto de errores:

\begin{itemize}[leftmargin=*]
    \item \textbf{Validación de ISBN:} Regex para ISBN-10 e ISBN-13
    \item \textbf{Validación de año:} Debe estar entre 1000 y año actual
    \item \textbf{Validación de campos vacíos:} Ningún campo puede estar vacío
    \item \textbf{Manejo de archivos:} \texttt{try-except-finally} para I/O
    \item \textbf{Duplicados:} Verificar ISBN duplicado antes de agregar
    \item \textbf{Excepciones personalizadas:} Crea al menos 2 excepciones custom
\end{itemize}

Ejemplo de excepción personalizada:
\begin{lstlisting}[caption=Excepciones personalizadas]
class ISBNInvalidoError(Exception):
    """Excepcion para ISBN invalido"""
    pass

class LibroDuplicadoError(Exception):
    """Excepcion cuando se intenta agregar libro duplicado"""
    pass

class LibroNoEncontradoError(Exception):
    """Excepcion cuando no se encuentra un libro"""
    pass
\end{lstlisting}

\subsection{Módulo 5: Archivos y Persistencia}

\begin{itemize}[leftmargin=*]
    \item \textbf{Archivo principal:} \texttt{biblioteca.csv} para guardar todos los libros
    \item \textbf{Auto-guardado:} Guardar automáticamente después de cada cambio
    \item \textbf{Carga al inicio:} Cargar datos existentes al iniciar la aplicación
    \item \textbf{Exportar:} Permitir exportar a CSV personalizado
    \item \textbf{Importar:} Permitir importar desde CSV con validación
    \item \textbf{Backup:} Crear copia de seguridad antes de importar
    \item \textbf{Formato CSV:} ISBN, Título, Autor, Género, Año, Prestado, Prestado\_a, Fecha\_Prestamo
\end{itemize}

\subsection{Módulo 6: Expresiones Regulares}

Usa regex para:

\begin{itemize}[leftmargin=*]
    \item \textbf{Validar ISBN:} 
    \begin{itemize}
        \item ISBN-10: \texttt{\^{}(?:\textbackslash d\{9\}X|\textbackslash d\{10\})\$}
        \item ISBN-13: \texttt{\^{}\textbackslash d\{13\}\$}
    \end{itemize}
    \item \textbf{Búsqueda flexible:} Permitir búsqueda parcial insensible a mayúsculas
    \item \textbf{Validar año:} Solo dígitos, 4 caracteres
    \item \textbf{Extraer información:} De strings complejos durante importación
\end{itemize}

\subsection{Módulo 7: Pruebas Unitarias}

Crea pruebas con \texttt{pytest} para:

\begin{itemize}[leftmargin=*]
    \item \textbf{Clase Libro:}
    \begin{itemize}
        \item Test de creación de libro válido
        \item Test de ISBN inválido (debe lanzar excepción)
        \item Test de año inválido
        \item Test de préstamo y devolución
        \item Test de métodos \texttt{\_\_str\_\_} y \texttt{\_\_repr\_\_}
    \end{itemize}
    
    \item \textbf{Clase Biblioteca:}
    \begin{itemize}
        \item Test de agregar libro
        \item Test de agregar libro duplicado (debe lanzar excepción)
        \item Test de búsqueda por título
        \item Test de búsqueda por autor con regex
        \item Test de eliminar libro
        \item Test de estadísticas
        \item Test de guardar y cargar desde archivo
    \end{itemize}
    
    \item \textbf{Validaciones:}
    \begin{itemize}
        \item Test de validación de ISBN (casos válidos e inválidos)
        \item Test de validación de año
        \item Test de campos vacíos
    \end{itemize}
\end{itemize}

Estructura de archivos de prueba:
\begin{lstlisting}[caption=Estructura de pruebas]
tests/
    __init__.py
    test_libro.py
    test_biblioteca.py
    test_validaciones.py
    conftest.py  # Fixtures compartidas
\end{lstlisting}

\section{Requisitos Técnicos Detallados}

\subsection{Estructura del Proyecto}

Tu proyecto debe tener la siguiente estructura:

\begin{lstlisting}[language=bash, caption=Estructura de directorios]
biblioteca_personal/
    main.py                 # Punto de entrada
    libro.py               # Clase Libro
    biblioteca.py          # Clase Biblioteca
    gui.py                 # Interfaz grafica
    validaciones.py        # Funciones de validacion
    excepciones.py         # Excepciones personalizadas
    utilidades.py          # Funciones auxiliares
    config.py              # Configuracion
    requirements.txt       # Dependencias
    README.md             # Documentacion
    tests/                # Pruebas unitarias
        __init__.py
        test_libro.py
        test_biblioteca.py
        test_validaciones.py
        conftest.py
    datos/                # Archivos de datos
        biblioteca.csv
        backup/
\end{lstlisting}

\subsection{Uso de Conceptos por Módulo}

\subsubsection{Variables y Tipos de Datos (Semana 0)}
\begin{itemize}[leftmargin=*]
    \item Usar f-strings para todo el formateo
    \item Variables descriptivas en todo el código
    \item Conversión de tipos cuando sea necesario
    \item Métodos de strings: \texttt{strip()}, \texttt{title()}, \texttt{upper()}, \texttt{lower()}
\end{itemize}

\subsubsection{Condicionales (Semana 1)}
\begin{itemize}[leftmargin=*]
    \item Validaciones complejas con múltiples condiciones
    \item Uso de operadores lógicos (\texttt{and}, \texttt{or}, \texttt{not})
    \item Match-case para menús o categorización (Python 3.10+)
\end{itemize}

\subsubsection{Bucles (Semana 2)}
\begin{itemize}[leftmargin=*]
    \item Iterar sobre colecciones de libros
    \item Validación de entrada con bucles \texttt{while}
    \item Uso de \texttt{break} y \texttt{continue}
    \item Bucles \texttt{for} con \texttt{range()} y sobre iterables
\end{itemize}

\subsubsection{Excepciones (Semana 3)}
\begin{itemize}[leftmargin=*]
    \item Bloques \texttt{try-except-else-finally} en I/O de archivos
    \item Excepciones personalizadas heredando de \texttt{Exception}
    \item Uso de \texttt{raise} para lanzar excepciones
    \item Manejo específico de excepciones (\texttt{FileNotFoundError}, \texttt{ValueError}, etc.)
\end{itemize}

\subsubsection{Librerías (Semana 4)}
\begin{itemize}[leftmargin=*]
    \item \texttt{tkinter}: Interfaz gráfica completa
    \item \texttt{csv}: Lectura y escritura de archivos
    \item \texttt{datetime}: Registro de fechas de préstamo
    \item \texttt{re}: Expresiones regulares para búsqueda y validación
    \item \texttt{statistics}: Cálculo de estadísticas
    \item \texttt{os}: Manejo de rutas y archivos
    \item \texttt{sys}: Información del sistema
\end{itemize}

\subsubsection{Pruebas Unitarias (Semana 5)}
\begin{itemize}[leftmargin=*]
    \item Mínimo 20 tests con \texttt{pytest}
    \item Uso de \texttt{assert} para verificaciones
    \item Fixtures con \texttt{@pytest.fixture}
    \item Tests de excepciones con \texttt{pytest.raises}
    \item Cobertura de código > 80\%
\end{itemize}

\subsubsection{Archivos (Semana 6)}
\begin{itemize}[leftmargin=*]
    \item Lectura y escritura de CSV con módulo \texttt{csv}
    \item Uso de \texttt{with} para manejo automático de archivos
    \item Creación de backups automáticos
    \item Manejo de rutas con \texttt{os.path}
\end{itemize}

\subsubsection{Expresiones Regulares (Semana 7)}
\begin{itemize}[leftmargin=*]
    \item \texttt{re.match()} para validar ISBN
    \item \texttt{re.search()} para búsquedas en títulos/autores
    \item \texttt{re.findall()} para extraer información
    \item Flags: \texttt{re.IGNORECASE}
\end{itemize}

\subsubsection{POO (Semana 8)}
\begin{itemize}[leftmargin=*]
    \item Mínimo 2 clases principales: \texttt{Libro} y \texttt{Biblioteca}
    \item Encapsulación con atributos privados (\_atributo)
    \item Properties con \texttt{@property}
    \item Métodos de clase con \texttt{@classmethod}
    \item Métodos estáticos con \texttt{@staticmethod}
    \item Métodos especiales: \texttt{\_\_init\_\_}, \texttt{\_\_str\_\_}, \texttt{\_\_repr\_\_}
    \item Bonus: Herencia (crear subclases de Libro: LibroFisico, LibroDigital)
\end{itemize}

\subsubsection{Conceptos Avanzados (Semana 9)}
\begin{itemize}[leftmargin=*]
    \item List comprehensions para filtrar y transformar listas
    \item Dict comprehensions para crear diccionarios
    \item Funciones lambda en \texttt{sort()}, \texttt{filter()}, \texttt{map()}
    \item Sets para operaciones de conjuntos (géneros únicos, etc.)
    \item \texttt{*args} y \texttt{**kwargs} en funciones flexibles
    \item Decorador personalizado (ej: \texttt{@log\_operacion})
\end{itemize}

\section{Ejemplo de Código: Integración de Conceptos}

\begin{lstlisting}[caption=Ejemplo de método que integra múltiples conceptos]
import re
import csv
from datetime import datetime
from typing import List, Dict, Optional

class Biblioteca:
    def buscar_libros_avanzado(
        self, 
        patron: str, 
        campo: str = "titulo",
        usar_regex: bool = True,
        **filtros
    ) -> List[Dict]:
        """
        Busqueda avanzada de libros.
        
        Integra: regex, excepciones, comprehensions, 
        *args/**kwargs, type hints
        """
        try:
            # Validar campo de busqueda
            campos_validos = {"titulo", "autor", "genero"}
            if campo not in campos_validos:
                raise ValueError(
                    f"Campo debe ser uno de: {campos_validos}"
                )
            
            # Preparar patron de busqueda
            if usar_regex:
                patron_compilado = re.compile(
                    patron, 
                    re.IGNORECASE
                )
            else:
                patron = patron.lower()
            
            # Buscar usando comprehension
            resultados = [
                libro for libro in self._libros.values()
                if self._coincide(
                    libro, 
                    campo, 
                    patron, 
                    patron_compilado if usar_regex else None
                )
            ]
            
            # Aplicar filtros adicionales
            if filtros:
                resultados = self._aplicar_filtros(
                    resultados, 
                    **filtros
                )
            
            # Ordenar usando lambda
            resultados.sort(
                key=lambda libro: getattr(libro, campo)
            )
            
            return [libro.to_dict() for libro in resultados]
            
        except re.error as e:
            raise ValueError(
                f"Expresion regular invalida: {e}"
            )
        except Exception as e:
            raise RuntimeError(
                f"Error en busqueda: {e}"
            )
\end{lstlisting}

\section{Interfaz Gráfica con Tkinter}

\subsection{Ventana Principal - Ejemplo}

\begin{lstlisting}[caption=Estructura básica de la interfaz]
import tkinter as tk
from tkinter import ttk, messagebox

class BibliotecaGUI:
    def __init__(self, biblioteca):
        self.biblioteca = biblioteca
        self.ventana = tk.Tk()
        self.ventana.title("Sistema de Biblioteca Personal")
        self.ventana.geometry("1000x600")
        
        self._crear_menu()
        self._crear_barra_herramientas()
        self._crear_tabla_libros()
        self._crear_panel_detalles()
        self._crear_barra_estado()
        
        self.actualizar_lista()
    
    def _crear_menu(self):
        menubar = tk.Menu(self.ventana)
        self.ventana.config(menu=menubar)
        
        # Menu Archivo
        menu_archivo = tk.Menu(menubar, tearoff=0)
        menubar.add_cascade(label="Archivo", menu=menu_archivo)
        menu_archivo.add_command(
            label="Importar CSV", 
            command=self.importar_csv
        )
        menu_archivo.add_command(
            label="Exportar CSV", 
            command=self.exportar_csv
        )
        menu_archivo.add_separator()
        menu_archivo.add_command(
            label="Salir", 
            command=self.ventana.quit
        )
        
        # Menu Gestion
        menu_gestion = tk.Menu(menubar, tearoff=0)
        menubar.add_cascade(label="Gestion", menu=menu_gestion)
        menu_gestion.add_command(
            label="Agregar Libro", 
            command=self.agregar_libro
        )
        menu_gestion.add_command(
            label="Editar Libro", 
            command=self.editar_libro
        )
        menu_gestion.add_command(
            label="Eliminar Libro", 
            command=self.eliminar_libro
        )
        
        # Menu Estadisticas
        menu_stats = tk.Menu(menubar, tearoff=0)
        menubar.add_cascade(
            label="Estadisticas", 
            menu=menu_stats
        )
        menu_stats.add_command(
            label="Ver Estadisticas", 
            command=self.mostrar_estadisticas
        )
    
    def _crear_tabla_libros(self):
        # Frame para la tabla
        frame_tabla = ttk.Frame(self.ventana)
        frame_tabla.pack(fill=tk.BOTH, expand=True, padx=10, pady=10)
        
        # Scrollbars
        scroll_y = ttk.Scrollbar(frame_tabla, orient=tk.VERTICAL)
        scroll_x = ttk.Scrollbar(frame_tabla, orient=tk.HORIZONTAL)
        
        # Treeview
        self.tabla = ttk.Treeview(
            frame_tabla,
            columns=("ISBN", "Titulo", "Autor", "Genero", "Anio"),
            show="headings",
            yscrollcommand=scroll_y.set,
            xscrollcommand=scroll_x.set
        )
        
        # Configurar columnas
        self.tabla.heading("ISBN", text="ISBN")
        self.tabla.heading("Titulo", text="Titulo")
        self.tabla.heading("Autor", text="Autor")
        self.tabla.heading("Genero", text="Genero")
        self.tabla.heading("Anio", text="Anio")
        
        self.tabla.column("ISBN", width=120)
        self.tabla.column("Titulo", width=300)
        self.tabla.column("Autor", width=200)
        self.tabla.column("Genero", width=100)
        self.tabla.column("Anio", width=80)
        
        # Empaquetar
        scroll_y.config(command=self.tabla.yview)
        scroll_x.config(command=self.tabla.xview)
        scroll_y.pack(side=tk.RIGHT, fill=tk.Y)
        scroll_x.pack(side=tk.BOTTOM, fill=tk.X)
        self.tabla.pack(fill=tk.BOTH, expand=True)
        
        # Doble click para editar
        self.tabla.bind("<Double-1>", lambda e: self.editar_libro())
\end{lstlisting}

\subsection{Validación en Tiempo Real}

\begin{lstlisting}[caption=Validación de ISBN en tiempo real]
def _validar_isbn_tiempo_real(self, event):
    """
    Valida ISBN mientras el usuario escribe.
    Integra: regex, tkinter, excepciones
    """
    isbn = self.entry_isbn.get().strip()
    
    if not isbn:
        self.label_isbn_estado.config(
            text="", 
            foreground="black"
        )
        return
    
    try:
        if Libro.validar_isbn(isbn):
            self.label_isbn_estado.config(
                text="ISBN valido", 
                foreground="green"
            )
            self.btn_guardar.config(state=tk.NORMAL)
        else:
            raise ISBNInvalidoError("Formato incorrecto")
    except ISBNInvalidoError as e:
        self.label_isbn_estado.config(
            text=f"ISBN invalido: {e}", 
            foreground="red"
        )
        self.btn_guardar.config(state=tk.DISABLED)
\end{lstlisting}

\section{Funcionalidades Adicionales Requeridas}

\subsection{Sistema de Préstamos}
\begin{itemize}[leftmargin=*]
    \item Registrar préstamo con nombre de persona y fecha
    \item Calcular días transcurridos desde el préstamo
    \item Alerta si un libro lleva más de 30 días prestado
    \item Historial de préstamos anteriores
\end{itemize}

\subsection{Estadísticas Avanzadas}
\begin{itemize}[leftmargin=*]
    \item Total de libros por género (usar \texttt{collections.Counter})
    \item Año promedio de publicación (usar \texttt{statistics.mean})
    \item Top 5 autores con más libros
    \item Porcentaje de libros prestados vs. disponibles
    \item Libro más antiguo y más reciente
\end{itemize}

\subsection{Búsqueda y Filtrado}
\begin{itemize}[leftmargin=*]
    \item Búsqueda simple (substring)
    \item Búsqueda con regex
    \item Filtro por rango de años
    \item Filtro por género
    \item Filtro por estado (prestado/disponible)
    \item Combinar múltiples filtros
\end{itemize}

\subsection{Importación/Exportación}
\begin{itemize}[leftmargin=*]
    \item Exportar toda la biblioteca a CSV
    \item Exportar solo libros filtrados
    \item Importar desde CSV con validación
    \item Crear backup automático antes de importar
    \item Reporte de errores durante importación
\end{itemize}

\section{Características de Calidad}

\subsection{Documentación}
\begin{itemize}[leftmargin=*]
    \item \textbf{Docstrings:} Todas las clases y funciones públicas deben tener docstrings
    \item \textbf{Type hints:} Usar anotaciones de tipo en todas las funciones
    \item \textbf{README.md:} Documentación completa del proyecto
    \item \textbf{Comentarios:} Explicar lógica compleja
\end{itemize}

\subsection{Código Limpio}
\begin{itemize}[leftmargin=*]
    \item Nombres descriptivos para variables, funciones y clases
    \item Funciones pequeñas y enfocadas (máx. 30-40 líneas)
    \item Evitar código duplicado (DRY - Don't Repeat Yourself)
    \item Indentación consistente (4 espacios)
    \item Líneas de máximo 79-100 caracteres
\end{itemize}

\subsection{Manejo de Errores}
\begin{itemize}[leftmargin=*]
    \item Nunca dejar bloques \texttt{except} vacíos
    \item Mensajes de error claros y útiles
    \item Logging de errores para debugging
    \item Validación de entrada del usuario
\end{itemize}

\section{README.md Requerido}

Tu archivo README.md debe incluir:

\begin{lstlisting}[language=bash, caption=Contenido del README]
# Sistema de Gestion de Biblioteca Personal

## Descripcion
[Descripcion breve del proyecto]

## Caracteristicas
- Lista de funcionalidades principales

## Requisitos
- Python 3.10+
- Librerias: tkinter, pytest, etc.

## Instalacion
```bash
pip install -r requirements.txt
```

## Uso
```bash
python main.py
```

## Estructura del Proyecto
[Descripcion de archivos y carpetas]

## Pruebas
```bash
pytest tests/ -v
pytest --cov=. tests/  # Con cobertura
```

## Conceptos de Python Utilizados
[Lista de todos los conceptos integrados]

## Autor
[Tu nombre]

## Fecha
[Fecha de entrega]
\end{lstlisting}

\section{Entregables}

\subsection{Código Fuente}
\begin{enumerate}[leftmargin=*]
    \item Todos los archivos .py mencionados en la estructura
    \item Archivo \texttt{requirements.txt} con dependencias
    \item Archivo \texttt{README.md} completo
    \item Carpeta \texttt{tests/} con todas las pruebas
    \item Archivo CSV de ejemplo con datos de prueba
\end{enumerate}

\subsection{Documentación}
\begin{enumerate}[leftmargin=*]
    \item README.md completo
    \item Docstrings en todas las clases y funciones públicas
    \item Comentarios explicativos en lógica compleja
    \item Manual de usuario (opcional, puntos extra)
\end{enumerate}

\subsection{Pruebas}
\begin{enumerate}[leftmargin=*]
    \item Mínimo 20 pruebas unitarias
    \item Cobertura de código > 80\%
    \item Todas las pruebas deben pasar
    \item Incluir tests de excepciones
\end{enumerate}

\section{Defensa del Proyecto (60\% de la nota final)}

\subsection{Formato de la Defensa}

La defensa es \textbf{obligatoria} y consta de:

\begin{enumerate}[leftmargin=*]
    \item \textbf{Presentación (10 minutos):}
    \begin{itemize}
        \item Demo completa de la aplicación funcionando
        \item Mostrar todas las funcionalidades implementadas
        \item Ejecutar suite de pruebas
    \end{itemize}
    
    \item \textbf{Explicación del código (15 minutos):}
    \begin{itemize}
        \item Arquitectura general del proyecto
        \item Explicación de clases principales
        \item Decisiones de diseño importantes
        \item Cómo integraste cada concepto del curso
    \end{itemize}
    
    \item \textbf{Preguntas técnicas (15 minutos):}
    \begin{itemize}
        \item Explicar partes específicas del código
        \item Modificar código en vivo
        \item Resolver bugs introducidos intencionalmente
        \item Agregar funcionalidad pequeña
    \end{itemize}
\end{enumerate}

\subsection{Preguntas Tipo de la Defensa}

Prepárate para responder preguntas como:

\begin{itemize}[leftmargin=*]
    \item ``Explica cómo funciona tu validación de ISBN''
    \item ``¿Qué pasa si cambias este \texttt{float()} por \texttt{int()}?''
    \item ``¿Por qué usaste un diccionario en lugar de una lista?''
    \item ``Muestra dónde usaste list comprehensions y explica por qué''
    \item ``¿Cómo manejas el caso de un archivo CSV corrupto?''
    \item ``Explica esta expresión regular''
    \item ``¿Qué hace este decorador?''
    \item ``Modifica esta búsqueda para que sea case-sensitive''
    \item ``Agrega un nuevo campo a la clase Libro''
\end{itemize}

\subsection{Demostración de Comprensión}

Durante la defensa, debes demostrar que:

\begin{itemize}[leftmargin=*]
    \item Entiendes cada línea de código que escribiste
    \item Puedes explicar por qué elegiste cada solución
    \item Conoces alternativas y sus trade-offs
    \item Puedes modificar el código en tiempo real
    \item Comprendes cómo los conceptos se integran
\end{itemize}

\section{Rúbrica de Evaluación}

\subsection{Distribución General}
\begin{itemize}
    \item \textbf{Implementación (40\%):} 40 puntos
    \item \textbf{Defensa (60\%):} 60 puntos
    \item \textbf{Total:} 100 puntos
    \item \textbf{Puntos Extra:} Hasta +20 puntos
\end{itemize}

\subsection{Parte 1: Implementación (40 puntos)}

\subsubsection{1.1 Funcionalidad Core (16 puntos)}
\begin{itemize}[leftmargin=*]
    \item \textbf{CRUD completo (8 pts):} Crear, leer, actualizar, eliminar libros funciona perfectamente
    \item \textbf{Persistencia (4 pts):} Guarda y carga datos correctamente
    \item \textbf{Búsqueda (4 pts):} Búsqueda funcional con regex
\end{itemize}

\subsubsection{1.2 Programación Orientada a Objetos (8 puntos)}
\begin{itemize}[leftmargin=*]
    \item \textbf{Diseño de clases (3 pts):} Clases bien diseñadas y cohesivas
    \item \textbf{Encapsulación (2 pts):} Uso correcto de atributos privados y properties
    \item \textbf{Métodos especiales (2 pts):} \texttt{\_\_init\_\_}, \texttt{\_\_str\_\_}, \texttt{\_\_repr\_\_}
    \item \textbf{Métodos de clase/estáticos (1 pt):} Al menos uno de cada tipo
\end{itemize}

\subsubsection{1.3 Manejo de Errores y Validaciones (6 puntos)}
\begin{itemize}[leftmargin=*]
    \item \textbf{Excepciones personalizadas (2 pts):} Al menos 2, bien implementadas
    \item \textbf{Validaciones robustas (3 pts):} ISBN, año, campos vacíos, etc.
    \item \textbf{Manejo de I/O (1 pt):} Try-except apropiado en archivos
\end{itemize}

\subsubsection{1.4 Interfaz Gráfica (6 puntos)}
\begin{itemize}[leftmargin=*]
    \item \textbf{Diseño funcional (3 pts):} Interfaz completa y usable
    \item \textbf{Validación en tiempo real (2 pts):} Feedback inmediato
    \item \textbf{UX (1 pt):} Fácil de usar, mensajes claros
\end{itemize}

\subsubsection{1.5 Calidad de Código (4 puntos)}
\begin{itemize}[leftmargin=*]
    \item \textbf{Documentación (2 pts):} Docstrings, comentarios, README
    \item \textbf{Organización (1 pt):} Estructura de archivos clara
    \item \textbf{Estilo (1 pt):} Código limpio y legible
\end{itemize}

\subsection{Parte 2: Defensa del Proyecto (60 puntos)}

\subsubsection{2.1 Demostración (18 puntos)}
\begin{itemize}[leftmargin=*]
    \item \textbf{Funcionalidad completa (10 pts):} Demuestra todas las funciones
    \item \textbf{Manejo de casos (5 pts):} Muestra casos normales y edge cases
    \item \textbf{Pruebas (3 pts):} Ejecuta suite de pruebas exitosamente
\end{itemize}

\subsubsection{2.2 Explicación Técnica (24 puntos)}
\begin{itemize}[leftmargin=*]
    \item \textbf{Arquitectura (6 pts):} Explica estructura general y flujo
    \item \textbf{POO (6 pts):} Explica diseño de clases y decisiones
    \item \textbf{Integración de conceptos (6 pts):} Muestra cómo usó cada tema del curso
    \item \textbf{Decisiones de diseño (6 pts):} Justifica elecciones técnicas
\end{itemize}

\subsubsection{2.3 Respuestas a Preguntas (18 puntos)}
\begin{itemize}[leftmargin=*]
    \item \textbf{Comprensión del código (10 pts):} Explica cualquier parte solicitada
    \item \textbf{Modificación en vivo (5 pts):} Puede cambiar código efectivamente
    \item \textbf{Resolución de problemas (3 pts):} Debugging y solución de bugs
\end{itemize}

\subsection{Parte 3: Pruebas Unitarias (Requisito, no puntuado separadamente)}

\textbf{Mínimo requerido:}
\begin{itemize}[leftmargin=*]
    \item 20 pruebas unitarias
    \item Cobertura > 80\%
    \item Todas las pruebas pasan
\end{itemize}

\textbf{Nota:} Si no cumples los requisitos mínimos de pruebas, se penaliza la sección de implementación.

\subsection{Puntos Extra (Hasta +20 puntos)}

\subsubsection{Funcionalidades Avanzadas (+10 pts)}
\begin{itemize}[leftmargin=*]
    \item \textbf{Herencia (+3 pts):} Subclases de Libro (LibroFisico, LibroDigital)
    \item \textbf{Gráficos (+3 pts):} Visualización de estadísticas con matplotlib
    \item \textbf{Base de datos (+4 pts):} Usar SQLite en lugar de CSV
\end{itemize}

\subsubsection{Calidad Excepcional (+10 pts)}
\begin{itemize}[leftmargin=*]
    \item \textbf{Documentación completa (+3 pts):} Manual de usuario detallado
    \item \textbf{Cobertura > 95\% (+2 pts):} Pruebas exhaustivas
    \item \textbf{Temas (+2 pts):} Interfaz con temas personalizables
    \item \textbf{Logging (+3 pts):} Sistema de logs completo
\end{itemize}

\subsection{Penalizaciones}

\begin{itemize}[leftmargin=*]
    \item \textbf{-100\% defensa:} No asiste a defensa (máx. 40/100)
    \item \textbf{-50\% defensa:} No puede explicar código (posible plagio)
    \item \textbf{-10 pts:} Menos de 20 pruebas unitarias
    \item \textbf{-10 pts:} Cobertura < 80\%
    \item \textbf{-5 pts:} Sin README.md o incompleto
    \item \textbf{-5 pts:} Estructura de proyecto incorrecta
    \item \textbf{-3 pts por día:} Entrega tardía (máx. 3 días)
\end{itemize}

\subsection{Nota Mínima Aprobatoria}

Para aprobar el proyecto final necesitas obtener al menos \textbf{70 puntos de 100}.

\subsection{Escala de Calificación}

\begin{itemize}
    \item \textbf{90-100+ pts:} Excelente - Proyecto profesional
    \item \textbf{80-89 pts:} Muy Bueno - Dominio sólido
    \item \textbf{70-79 pts:} Bueno - Competente
    \item \textbf{60-69 pts:} Suficiente - Necesita mejorar
    \item \textbf{0-59 pts:} Insuficiente - No aprobado
\end{itemize}

\section{Cronograma Sugerido}

\subsection{Semana 1}
\begin{itemize}[leftmargin=*]
    \item Diseñar arquitectura del proyecto
    \item Implementar clases Libro y Biblioteca
    \item Crear funciones de validación
    \item Escribir primeras pruebas unitarias
\end{itemize}

\subsection{Semana 2}
\begin{itemize}[leftmargin=*]
    \item Implementar interfaz gráfica básica
    \item Agregar funcionalidad CRUD
    \item Implementar persistencia con CSV
    \item Completar pruebas unitarias
    \item Escribir documentación
    \item Preparar defensa
\end{itemize}

\section{Recursos de Apoyo}

\subsection{Documentación Oficial}
\begin{itemize}[leftmargin=*]
    \item Python: \url{https://docs.python.org/3/}
    \item Tkinter: \url{https://docs.python.org/3/library/tkinter.html}
    \item Pytest: \url{https://docs.pytest.org/}
    \item CSV: \url{https://docs.python.org/3/library/csv.html}
    \item Regex: \url{https://docs.python.org/3/library/re.html}
\end{itemize}

\subsection{Tutoriales Recomendados}
\begin{itemize}[leftmargin=*]
    \item Real Python - Tkinter GUI Programming
    \item Corey Schafer - OOP en Python (YouTube)
    \item Python Testing with pytest (Brian Okken)
\end{itemize}

\section{Preguntas Frecuentes}

\textbf{P: ¿Puedo usar una base de datos en lugar de CSV?}\\
R: Sí, pero CSV es requisito mínimo. Si usas BD, también debes mantener export/import CSV y obtienes puntos extra.

\textbf{P: ¿Cuántas líneas de código debería tener?}\\
R: No hay mínimo/máximo, pero típicamente 800-1500 líneas para un proyecto completo (incluyendo tests).

\textbf{P: ¿Puedo trabajar en equipo?}\\
R: No, este es un proyecto individual. Debes poder explicar cada línea.

\textbf{P: ¿Qué pasa si no termino todo?}\\
R: Entrega lo que tengas. Es mejor entregar algo funcional parcial que nada. La defensa determinará tu comprensión.

\textbf{P: ¿Puedo usar librerías adicionales?}\\
R: Sí, pero las principales funcionalidades deben usar las librerías del curso. Librerías adicionales son bonus.

\textbf{P: ¿Cómo sé si mi interfaz es suficientemente buena?}\\
R: Debe ser funcional, no necesariamente bonita. Si un usuario puede realizar todas las operaciones sin confusión, es suficiente.

\section{Consejos Finales}

\begin{itemize}[leftmargin=*]
    \item \textbf{Empieza temprano:} Este proyecto es complejo y toma tiempo.
    
    \item \textbf{Commits frecuentes:} Usa Git para versionar tu código.
    
    \item \textbf{Prueba constantemente:} No esperes al final para probar.
    
    \item \textbf{Documenta mientras programas:} Es más fácil que hacerlo al final.
    
    \item \textbf{Pide ayuda:} Si te trabas, consulta con el profesor o compañeros (pero escribe tu propio código).
    
    \item \textbf{Prioriza:} Primero implementa lo básico, luego añade extras.
    
    \item \textbf{Practica tu defensa:} Explica tu código en voz alta antes de la reunión.
    
    \item \textbf{Ten backup:} Guarda tu código en la nube (GitHub, Drive, etc.).
    
    \item \textbf{Lee los errores:} Python te dice exactamente qué está mal.
    
    \item \textbf{Diviértete:} Este es tu proyecto, hazlo a tu manera dentro de los requisitos.
\end{itemize}

\vspace{1cm}

\begin{center}
\textbf{¡Éxito en tu proyecto final!}\\
\vspace{0.5cm}
\textit{Este proyecto demuestra todo lo que has aprendido.\\
Es tu oportunidad de brillar y mostrar tus habilidades como programador.\\
Confía en ti mismo y en todo lo que has aprendido durante el curso.}
\vspace{0.5cm}

\Large{\textbf{¡Adelante!}}
\end{center}

\end{document}
